\documentclass[11pt]{article}
\usepackage{varnernotes}

\newcommand{\R}{\mathbb{R}}
\newcommand{\E}{\mathbb{E}}
\newcommand{\Cov}{\operatorname{Cov}}
\newcommand{\Var}{\operatorname{Var}}
\newcommand{\np}{|\mathcal P|}

\begin{document}

\lecturenotes{6a}{Introduction to Single-Index Models}{Fall 2026}{September 29, 2026}{Jeffrey D. Varner}

\learningobjectives
  {Decompose asset growth rates}
  {Write the single-index model and its assumptions, interpret alpha, beta, and the residual, and split an asset's growth-rate variance into a market part and an asset-specific part.}
  {Build the SIM covariance}
  {Derive the rank-one-plus-diagonal mean and covariance structure implied by a single-index model, count what it saves, and state the residual assumption that must be checked.}
  {Estimate and diagnose a SIM}
  {Fit alpha and beta by least squares, compute their classical uncertainty in growth-rate units, and use residuals, confidence intervals, and $R^2$ without overstating what they establish.}

\section{Why Replace a Full Covariance Matrix?}

A portfolio $\mathcal P$ of $\np$ risky assets requires $\np$ expected growth rates and $\np(\np+1)/2$ distinct growth-rate covariance entries. The number of pairwise terms grows quadratically, estimates become noisy when the history is short relative to the asset universe, and an inverse covariance matrix can magnify that noise. Yet many securities visibly move together when market-wide news arrives. A factor model makes that shared source of movement explicit.

The single-index model (SIM), introduced by William Sharpe as a simplified model for portfolio analysis, compresses co-movement into exposure to one observed index. It is both a statistical regression and a covariance model. Those roles are related but not identical: fitting separate regressions does not by itself prove that their residuals are mutually uncorrelated.

\begin{figure}[ht]
  \centering
  \begin{tikzpicture}[>=Latex,font=\sffamily\small,node distance=1.05cm and 1.45cm]
    \node[draw=vnink,fill=vnsoft,minimum width=2.15cm,minimum height=0.78cm] (market) {market growth $g_M$};
    \node[draw=vnrule,fill=vnpanel,right=2.25cm of market,yshift=1.25cm,minimum width=2.15cm] (a1) {asset $1$: $g_1$};
    \node[draw=vnrule,fill=vnpanel,right=2.25cm of market,minimum width=2.15cm] (a2) {asset $2$: $g_2$};
    \node[draw=vnrule,fill=vnpanel,right=2.25cm of market,yshift=-1.25cm,minimum width=2.15cm] (am) {last asset: $g_{\np}$};
    \draw[vnlink,thick,->] (market.east) -- node[above,sloped] {$\beta_1$} (a1.west);
    \draw[vnlink,thick,->] (market.east) -- node[above] {$\beta_2$} (a2.west);
    \draw[vnlink,thick,->] (market.east) -- node[below,sloped] {$\beta_{\np}$} (am.west);
    \node[left=0.85cm of a1,text=vnmuted] (e1) {$\varepsilon_1$};
    \node[left=0.85cm of a2,text=vnmuted] (e2) {$\varepsilon_2$};
    \node[left=0.85cm of am,text=vnmuted] (em) {$\varepsilon_{\np}$};
    \draw[vnmuted,->] (e1)--(a1);
    \draw[vnmuted,->] (e2)--(a2);
    \draw[vnmuted,->] (em)--(am);
  \end{tikzpicture}
  \caption{One common market growth rate reaches every asset through its beta loading, while each residual records variation not explained by that index. The diagonal-residual assumption makes the residual arrows mutually uncorrelated; it should be checked rather than silently presumed.}
  \label{fig:sim-graph}
\end{figure}

\section{The Single-Index Growth Model}

Fix one sampling interval $\Delta t$, such as one trading day expressed in years. Let $g_M\in\R$ be the continuously compounded growth rate of a market index $M$ (SPY in this course; the letter names the index, and $\np$ counts the assets) and let $g_i\in\R$ be the growth rate of risky asset $i\in\mathcal P$, both on the same day. A SIM represents each asset as
\begin{equation}
  \label{eq:sim-scalar}
  g_i=\alpha_i+\beta_i g_M+\varepsilon_i.
\end{equation}
Here $\alpha_i$ and $\varepsilon_i$ are in growth-rate units, time$^{-1}$, while $\beta_i$ is a dimensionless market loading. The conventional population assumptions are
\begin{equation}
  \label{eq:sim-assumptions}
  \E[\varepsilon_i]=0,
  \qquad
  \Cov(\varepsilon_i,g_M)=0,
  \qquad
  \Cov(\varepsilon_i,\varepsilon_j)=0\quad(i\neq j).
\end{equation}
The first two conditions define a linear projection when the parameters are population least-squares coefficients. The third is an additional covariance-model restriction. Market proxies, sector effects, omitted factors, and synchronized shocks can leave correlated residuals and violate it.

\begin{definition}{Systematic and idiosyncratic components}{risk-components}
For the SIM in \eqnref{sim-scalar}, $\beta_i g_M$ is the asset's \emph{systematic component}: its variation is inherited from the common index. The residual $\varepsilon_i$ is the \emph{idiosyncratic component}: variation orthogonal to the chosen index under \eqnref{sim-assumptions}. Idiosyncratic does not mean predictable, harmless, or independent through time.
\end{definition}

If $\sigma_{g,M}^2:=\Var(g_M)>0$, the population linear-projection coefficients satisfy
\begin{equation}
  \label{eq:projection-coefficients}
  \beta_i=\frac{\Cov(g_i,g_M)}{\sigma_{g,M}^2}=\rho_{iM}\frac{\sigma_{g,i}}{\sigma_{g,M}},
  \qquad
  \alpha_i=\E[g_i]-\beta_i\E[g_M],
\end{equation}
with $\rho_{iM}$ the correlation of the two growth rates and $\sigma_{g,i}$ the asset's growth-rate standard deviation (positive). Taking the variance of \eqnref{sim-scalar} and using only $\Cov(\varepsilon_i,g_M)=0$ splits the asset's risk, $\Var(g_i)=\beta_i^2\sigma_{g,M}^2+\sigma_{g,\varepsilon,i}^2$ with $\sigma_{g,\varepsilon,i}^2:=\Var(\varepsilon_i)$: a systematic part inherited from the market and an idiosyncratic part, whose relative sizes are unrelated to the size of beta; the systematic fraction is $\rho_{iM}^2$. Thus beta measures sensitivity to the selected index and sampling convention, not total risk. A beta of $1.4$ does not mean the asset always moves by $1.4$ times the market; \eqnref{sim-scalar} includes an intercept and a residual on every observation. For the log return $r=(\Delta t)g$, the return-form coefficients satisfy $\alpha_{r,i}=(\Delta t)\alpha_i$, $\varepsilon_{r,i}=(\Delta t)\varepsilon_i$, and the same $\beta_i$.

\section{Mean and Covariance Implied by the SIM}

Let $\boldsymbol\alpha,\boldsymbol\beta\in\R^{\np}$ collect the intercepts and betas, let $\boldsymbol\varepsilon\in\R^{\np}$ collect the residuals, and let $D_g:=\operatorname{diag}(\sigma_{g,\varepsilon,1}^2,\ldots,\sigma_{g,\varepsilon,\np}^2)$. The vector model is
\begin{equation}
  \label{eq:sim-vector}
  \mathbf g=\boldsymbol\alpha+\boldsymbol\beta g_M+\boldsymbol\varepsilon.
\end{equation}
Under \eqnref{sim-assumptions}, its moments have the low-rank-plus-diagonal form in \propref{sim-moments}.

\begin{proposition}{Single-index moments}{sim-moments}
Let $\mathbf g\in\R^{\np}$ satisfy \eqnref{sim-vector}. Suppose $\E[\boldsymbol\varepsilon]=\mathbf0$, $\Cov(\boldsymbol\varepsilon,g_M)=\mathbf0$, and $\Cov(\boldsymbol\varepsilon)=D_g$. If $\bar g_M:=\E[g_M]$ and $\sigma_{g,M}^2:=\Var(g_M)$, then
\begin{equation}
  \label{eq:sim-moments}
  \bar{\mathbf g}:=\E[\mathbf g]
  =\boldsymbol\alpha+\boldsymbol\beta\bar g_M,
  \qquad
  \Sigma_g:=\Cov(\mathbf g)
  =\sigma_{g,M}^2\boldsymbol\beta\boldsymbol\beta^{\mathsf T}+D_g.
\end{equation}
Consequently, $\Cov(g_i,g_j)=\beta_i\beta_j\sigma_{g,M}^2$ for $i\neq j$, while $\Var(g_i)=\beta_i^2\sigma_{g,M}^2+\sigma_{g,\varepsilon,i}^2$.
\end{proposition}

The matrix $\sigma_{g,M}^2\boldsymbol\beta\boldsymbol\beta^{\mathsf T}$ has rank at most one. Instead of estimating every cross-covariance separately, the SIM estimates $\np$ betas, $\np$ residual growth variances, and one market growth variance ($2\np+1$ numbers for the covariance; $\np$ intercepts and the market mean more for the mean vector). This reduction buys stability only if the omitted residual correlations are small enough for the intended use. Fitting the $\np$ regressions separately makes every sample residual orthogonal to the market column and does nothing to make residuals mutually uncorrelated; in the course dataset the median absolute residual correlation is small, but related firms keep a few tenths and near-duplicate securities are almost perfectly correlated.

\section{Portfolio Risk Becomes Interpretable}

Let $\mathbf w\in\R^{\np}$ be fixed portfolio weights. The linearized one-period portfolio growth rate $g_p:=\mathbf w^{\mathsf T}\mathbf g$ (L5b) has market beta $\beta_p:=\mathbf w^{\mathsf T}\boldsymbol\beta$ and residual $\varepsilon_p:=\mathbf w^{\mathsf T}\boldsymbol\varepsilon$. Substituting \eqnref{sim-vector} gives
\begin{equation}
  \label{eq:portfolio-sim}
  g_p=\mathbf w^{\mathsf T}\boldsymbol\alpha+\beta_p g_M+\varepsilon_p,
  \qquad
  \Var(g_p)=\beta_p^2\sigma_{g,M}^2+\mathbf w^{\mathsf T}D_g\mathbf w.
\end{equation}
The first variance term is common-factor risk. The second is residual risk. For an equal-weight portfolio with $w_i=1/\np$ and residual variances bounded by a constant $K$, the residual variance satisfies
\begin{equation}
  \label{eq:diversification-bound}
  \mathbf w^{\mathsf T}D_g\mathbf w
  =\frac{1}{\np^2}\sum_{i\in\mathcal P}\sigma_{g,\varepsilon,i}^2
  \leq\frac{K}{\np}.
\end{equation}
Thus firm-specific variance can vanish as the portfolio broadens under the diagonal-residual assumption, while systematic variance generally remains. Concentrated weights or correlated residuals (a general residual covariance $\Omega_\varepsilon$ replaces $D_g$ and the bound fails) destroy this conclusion.

\section{Estimating Alpha and Beta}

For one asset, observe $N$ paired growth rates $\{(g_{M,t},g_{i,t})\}_{t=1}^N$ (the columns of the L5a data matrix). Let $\mathbf y\in\R^N$ contain the asset growth rates and let $X\in\R^{N\times2}$ have row $t$ equal to $(1,g_{M,t})$. When $X$ has full column rank, ordinary least squares estimates $\boldsymbol\theta_i:=(\alpha_i,\beta_i)^{\mathsf T}$ by
\begin{equation}
  \label{eq:ols}
  \widehat{\boldsymbol\theta}_i
  =(X^{\mathsf T}X)^{-1}X^{\mathsf T}\mathbf y,
  \qquad
  \widehat{\boldsymbol\varepsilon}_{i}
  =\mathbf y-X\widehat{\boldsymbol\theta}_i.
\end{equation}
Numerically, a QR or singular-value decomposition is preferable to explicitly forming the inverse in \eqnref{ols}; the three agree whenever $X$ has full column rank. The growth-residual variance estimate is $s_{g,\varepsilon,i}^2=\|\widehat{\boldsymbol\varepsilon}_{i}\|_2^2/(N-2)$ under the usual homoskedastic regression convention. Do not divide this value by $\Delta t$: the observations are already growth rates. The corresponding log-return residual variance is $(\Delta t)^2s_{g,\varepsilon,i}^2$, while its GBM covariance-rate contribution is $\Delta t\,s_{g,\varepsilon,i}^2$.

Under the classical assumptions (design held fixed, $\E[\boldsymbol\varepsilon\mid X]=\mathbf0$, independent homoskedastic Gaussian errors), $\widehat{\boldsymbol\theta}_i$ is unbiased with covariance $\sigma_{g,\varepsilon,i}^2(X^{\mathsf T}X)^{-1}$, estimated by $s_{g,\varepsilon,i}^2(X^{\mathsf T}X)^{-1}$; the standard errors are the square roots of its diagonal, and $\hat\theta_j\pm t_{0.975,N-2}\,\mathrm{SE}(\hat\theta_j)$ is a $95\%$ interval. Financial growth rates violate the list: they are heavy-tailed, their absolute values cluster in time, and growth rates built from volume-weighted average prices carry a positive autocorrelation at lag one. Heteroskedasticity/autocorrelation-consistent errors, block bootstrap intervals, or rolling estimates may then be more credible than the i.i.d. formula, and residual and parametric bootstraps (both independent-draws procedures) are a first check. Whatever method is used, the training window and sampling interval must be stated.

\notebooklink
  {L6a: Estimating Single Index Models from Historical Data}
  {https://github.com/varnerlab/CHEME-5660-CourseRepository-Fall-2026/blob/main/lectures/week-6/L6a/CHEME-5660-L6a-Example-SVD-SIM-Estimation-Fall-2026.ipynb}
  {Fit the SIM for one firm by least squares three equivalent ways, compute its classical uncertainty and $R^2$ next to those of an index fund, fit every security in the dataset, check the residual-correlation assumption, and save the parameter archive that L6b loads.}

\notebooklink
  {L6a: Bootstrap Uncertainty in a Single Index Model}
  {https://github.com/varnerlab/CHEME-5660-CourseRepository-Fall-2026/blob/main/lectures/week-6/L6a/CHEME-5660-L6a-Example-SIM-Parameter-Uncertainty-Fall-2026.ipynb}
  {Compare empirical-residual and Gaussian parametric bootstrap intervals with the classical standard errors, then audit the fitted residuals for heavy tails and dependence.}

\section{What \texorpdfstring{$R^2$}{R-squared}, Alpha, and Beta Do Not Prove}

In a least-squares regression with an intercept, $R_i^2:=1-\sum_t\widehat\varepsilon_{i,t}^2/\sum_t(g_{i,t}-g_i')^2$, with $g_i'$ the sample mean, is the sample fraction of asset-growth variation explained by the fitted market regressor, and equals the squared sample correlation $\hat\rho_{iM}^2$. It is not the probability that the model is true, a forecast-accuracy guarantee, or a measure of investment quality. A low $R^2$ can coexist with a useful beta estimate, and a high $R^2$ can coexist with unstable parameters or omitted nonlinear structure.

Likewise, a positive estimated alpha is an intercept conditional on the chosen market proxy, return convention, sample window, price series (volume-weighted average prices, not dividend-adjusted total returns), and specification. It is not automatically abnormal performance. Beta is exposure to the selected factor, not total risk. Before using a SIM covariance in an optimizer, inspect residual cross-correlations, rolling stability, out-of-sample error, and sensitivity to the market proxy.

The notebooks work with continuously compounded growth rates $g=r/\Delta t$. Means, intercepts, residuals, risk-free adjustments, and covariances must all use this convention. In particular, $\Sigma_r=(\Delta t)^2\Sigma_g$ and the GBM covariance rate is $C=\Delta t\,\Sigma_g$. Simple returns remain distinct from both quantities.

\newpage
\section*{Optional Advanced Material}

\notebooklink
  {L6a Advanced: Single Index Model Estimation Theory}
  {https://github.com/varnerlab/CHEME-5660-CourseRepository-Fall-2026/blob/main/lectures/week-6/L6a/advanced/sim/CHEME-5660-L6a-Advanced-SIM-Theory-Fall-2026.ipynb}
  {Derive the ridge estimator and its covariance, distinguish the residual and parametric bootstraps, keep the growth-rate, log-return, and volatility conventions straight, propagate parameter uncertainty into portfolio risk and weights, and write the maximum-Sharpe problem as a second-order cone program.}

\newpage
\thispagestyle{fancy}
\keytakeaways
  {In this lecture, we replaced a dense covariance description with one common market factor plus asset-specific residuals, derived the covariance and portfolio risk it implies, and connected that structure to least-squares estimation with its classical uncertainty.}
  {One factor splits risk in two}
  {An asset's growth-rate variance is $\beta_i^2\sigma_{g,M}^2+\sigma_{g,\varepsilon,i}^2$: beta measures systematic exposure, $R^2=\rho_{iM}^2$ measures how much the market explains, and the two can disagree.}
  {SIM is a covariance restriction}
  {The formula $\Sigma_g=\sigma_{g,M}^2\boldsymbol\beta\boldsymbol\beta^{\mathsf T}+D_g$ requires residuals to be uncorrelated across assets, not merely fitted separately; equal-weight portfolios then diversify residual risk and keep $\beta_p^2\sigma_{g,M}^2$.}
  {Least squares estimates, daily data qualify}
  {Standard errors and intervals are in growth-rate units with no $\Delta t$; heavy tails, lag-one dependence from volume-weighted prices, and volatility clustering make bootstrap or robust intervals the credible complement.}
  {Next, use the SIM moments inside risky-only and risky/risk-free minimum-variance allocation problems.}

\begin{readinglist}
  \item William F. Sharpe, \href{https://doi.org/10.1287/mnsc.9.2.277}{``A Simplified Model for Portfolio Analysis,''} \emph{Management Science} 9(2), 277--293 (1963).
  \item Edwin J. Elton, Martin J. Gruber, Stephen J. Brown, and William N. Goetzmann, \emph{Modern Portfolio Theory and Investment Analysis}, 9th ed. (2014), chapters on index models.
  \item John Y. Campbell, Andrew W. Lo, and A. Craig MacKinlay, \emph{The Econometrics of Financial Markets} (1997), chapters on return predictability and regression diagnostics.
\end{readinglist}

\vfill
{\sffamily\footnotesize\color{vnmuted}\textbf{Educational use and risk notice.} These notes are for instruction only and do not constitute investment advice. Financial decisions involve risk, and model outputs depend on assumptions.}

\end{document}
